\documentclass[11pt,a4paper]{letter}
\usepackage[margin=2.5cm]{geometry}
\usepackage[T1]{fontenc}
\usepackage[utf8]{inputenc}
\usepackage{hyperref}
\hypersetup{colorlinks=true,urlcolor=blue}

\begin{document}

\begin{letter}{The Editors \\
\textit{Journal of Translational Medicine}}

\opening{Dear Editors,}

We are pleased to submit our manuscript, ``\textbf{Limited plasma proteomic coverage constrains genetic drug-target prioritization in colorectal cancer}'', for consideration as a Research article in \textit{Journal of Translational Medicine}.

\textbf{What we did.} We built an end-to-end post-GWAS target-prioritization pipeline for colorectal cancer (CRC) that links blood cis-eQTL evidence (eQTLGen, $n = 31{,}684$) to a 78,473-case CRC GWAS, and then attempts to validate the resulting targets at the protein layer using two independent plasma proteomic panels (deCODE SomaScan and UKB-PPP Olink, about 35,000 participants each). We added tissue-level context (single-cell and Visium spatial transcriptomics of CRC), drug-tractability annotation, and PheWAS safety screening.

\textbf{What we found.} Of 15,582 genes tested, 75 passed Bonferroni-corrected SMR and 19 reached Tier~1 status, and 43 of 76 probe--gene pairs (57\%) colocalized at the eQTL layer. At the protein layer, however, the target list could barely be tested: of the 62 protein-coding genes, only 15 had matching plasma protein measurements (17 in the analysis set once two TNF-pathway genes are included), only \textbf{nine} had genome-wide significant cis-pQTL instruments available for Mendelian randomization, and only \textbf{one} (CCM2, PPH4 $= 0.9514$) colocalized with the CRC GWAS signal. Five of the six core non-MHC Tier~1 candidates have no protein measurement in either panel at all. The bottleneck for translating CRC GWAS loci into drug targets is therefore protein-layer coverage, not genetic significance.

\textbf{Why we believe this fits \textit{Journal of Translational Medicine}.} Our message is directly translational: we show that a widely used class of target-prioritization pipelines is limited by the resources used to validate it, and we turn that observation into a concrete reporting recommendation---report cis-pQTL coverage alongside any eQTL-derived target list, so that readers can see how much of the list was actually testable. We also provide the positive residue that survives this constraint: CCM2, the only gene with convergent eQTL fine-mapping, pQTL Mendelian randomization and pQTL colocalization evidence, and BMP2, which combines genetic support with the narrowest organ-system risk spectrum and a pre-existing pharmacological basis.

\textbf{Why the analysis is defensible.} All instrument sets were checked for strength (minimum $F = 18.2$, median $F$ 30.4--326.3), and we report heterogeneity and pleiotropy diagnostics rather than hiding them: Cochran's $Q$ indicates substantial heterogeneity for all five genes with multi-instrument estimates, and the MR-Egger intercept is non-significant for CCM2 ($p = 0.55$) and LIMA1 ($p = 0.11$) but significant for STAT6, CDKN1A and TNF. While preparing this manuscript we noticed that the MR-PRESSO outlier test is not resolvable at its default \texttt{NbDistribution = 1000} for instrument sets of this size (resolution $n/1000 = 0.061$--$0.228$, above the 0.05 threshold); we therefore re-ran it at \texttt{NbDistribution = 5000} (resolution at most 0.046) and report outlier counts together with the distortion test. We also report a power calculation explaining our FinnGen replication rate, and we report the 0-of-6 TCGA survival result rather than omitting it.

\textbf{Position relative to prior work.} We deliberately frame this study as an extension of, not a replacement for, existing CRC TWAS and colocalization work. Our genetic-layer findings overlap substantially with prior reports (44.9\% with Chen et~al., 2024; 17.4\% with Hazelwood et~al., 2025), and we say so explicitly. Our contribution is the quantification of the transcript-to-protein validation gap, not a new list of associated genes.

\textbf{Declarations.} This manuscript has not been published and is not under consideration elsewhere. All authors have approved the manuscript and agree with its submission to \textit{Journal of Translational Medicine}. Data availability and code availability statements are included. The authors declare no competing interests.

We would be glad to suggest reviewers with expertise in CRC genetics, statistical genetics and proteogenomic drug-target discovery; names can be supplied on request, or we are happy to leave the choice entirely to the editorial office.

Thank you for considering our manuscript.

\closing{Sincerely,}

\vspace{1cm}
Jiajie Zhou, MD \\
The Affiliated Huaian No.1 People's Hospital of Nanjing Medical University, Huaian, Jiangsu, China \\[0.4em]
E-mail: zjjwcwk2025@njmu.edu.cn; Tel: +86-15805233415 \\[0.4em]
On behalf of all co-authors

\end{letter}
\end{document}